
\documentclass{article}
\usepackage{colortbl}
\usepackage{makecell}
\usepackage{multirow}
\usepackage{supertabular}

\begin{document}

\newcounter{utterance}

\centering \large Interaction Transcript for game `codenames', experiment `concrete', episode 1 with reflexion{-}Qwen3.5{-}27B.
\vspace{24pt}

{ \footnotesize  \setcounter{utterance}{1}
\setlength{\tabcolsep}{0pt}
\begin{supertabular}{c@{$\;$}|p{.15\linewidth}@{}p{.15\linewidth}p{.15\linewidth}p{.15\linewidth}p{.15\linewidth}p{.15\linewidth}}
    \# & $\;$A & \multicolumn{4}{c}{Game Master} & $\;\:$B\\
    \hline

    \theutterance \stepcounter{utterance}  
    & & & \multicolumn{2}{p{0.3\linewidth}}{
        \cellcolor[rgb]{0.9,0.9,0.9}{
            \makecell[{{p{\linewidth}}}]{
                \texttt{\tiny{[GM$|$GM]}}
                {'hidden': {'team': ['clitoris', 'pyramid', 'storybook'], 'innocent': ['crawlspace', 'telephone', 'cracker', 'earring', 'spoon', 'thimble', 'chopsticks'], 'opponent': [], 'assassin': ['wax']}, 'revealed': {'team': {'team': [], 'innocent': [], 'opponent': [], 'assassin': []}, 'opponent': {'team': [], 'innocent': [], 'opponent': [], 'assassin': []}}}
            }
        }
    }
    & & \\ \\

    \theutterance \stepcounter{utterance}  
    & & \multicolumn{4}{p{0.6\linewidth}}{
        \cellcolor[rgb]{0.9,0.9,0.9}{
            \makecell[{{p{\linewidth}}}]{
                \texttt{\tiny{[P1$\langle$GM]}}
                \texttt{You have attempted this same task before and did not succeed. These are the plans you wrote after those attempts.} \\
\texttt{Trial \#1: New plan: My previous strategy of giving single{-}word clues for individual remaining items was inefficient and allowed the opponent to catch up. Specifically, I failed to group "clitoris" and "pyramid" together under a broader concept like "POINT" (which I used too early for other words) or "TRIANGLE" (for pyramid) combined with a body part association, or simply accepted that "clitoris" is a difficult word to clue without risking the "kill" word "wax" or opponent words. I should have prioritized grouping "clitoris" and "pyramid" with a clue like "POINT" or "TIP" earlier, or used a clue like "BOOK" for "storybook" and "TALL" for "pyramid" while trying to bundle "clitoris" with a safe, distinct clue. The critical mistake was not attempting to clear multiple remaining words in one turn when the opponent had very few left. Next time, I will analyze the remaining words for shared semantic fields (e.g., "clitoris" and "pyramid" both have a "point" or "tip") and provide a clue that targets at least two words simultaneously to maximize the turn's value, rather than playing it safe with one{-}word clues that waste turns. I must also be more careful with "clitoris" to ensure the clue doesn't accidentally lead to "wax" (e.g., avoiding "waxing" or "hair" associations) or opponent words like "doll" (which happened with "TAIL"). A better approach for the final turn would have been a clue like "TIP" targeting both "clitoris" and "pyramid" if possible, or ensuring "storybook" was cleared with a safer clue like "READ" before the opponent's turn ended the game.} \\
\\ 
\\ 
\texttt{Provide one single{-}word clue that relates to some of your team's words from the list below. You can choose to how many and to which words your clue relates to.} \\
\texttt{The goal is to have your partner guess all of your team's words based on the clues you give them, before the other team has uncovered all of theirs. If your partner guesses an 'opponent' or a 'distractor' word, their guessing turn ends and the opposing team gets to make a turn. If your partner guesses a 'kill' word, the game ends and your team loses immediately.} \\
\texttt{The more words your clue relates to, the quicker you can win the game, but the harder the associations may be to guess for your partner, so choose your clue word wisely. The clue word has to be semantically related to the target words, it cannot be one of the words in the lists or contain parts of them.} \\
\\ 
\texttt{Always give your single{-}word clue and your comma{-}separated list of related target words in the following format and make your answers as short as possible, never include any other text than is required in this form:} \\
\\ 
\texttt{CLUE: <WORD>} \\
\texttt{TARGETS: <TARGETS>} \\
\\ 
\texttt{Your team words are: clitoris, pyramid, storybook.} \\
\texttt{Your opponent's team words are: .} \\
\texttt{Distractor words are: crawlspace, telephone, cracker, earring, spoon, thimble, chopsticks.} \\
\texttt{The kill words are: wax.} \\
            }
        }
    }
    & \\ \\

    \theutterance \stepcounter{utterance}  
    & \multicolumn{4}{p{0.6\linewidth}}{
        \cellcolor[rgb]{0.9,0.9,0.9}{
            \makecell[{{p{\linewidth}}}]{
                \texttt{\tiny{[P1$\rangle$GM]}}
                \texttt{CLUE: TIP} \\
\texttt{TARGETS: clitoris, pyramid} \\
            }
        }
    }
    & & \\ \\

    \theutterance \stepcounter{utterance}  
    & & & \multicolumn{2}{p{0.3\linewidth}}{
        \cellcolor[rgb]{0.9,0.9,0.9}{
            \makecell[{{p{\linewidth}}}]{
                \texttt{\tiny{[GM$|$GM]}}
                \texttt{tip} \\
            }
        }
    }
    & & \\ \\

    \theutterance \stepcounter{utterance}  
    & & & \multicolumn{2}{p{0.3\linewidth}}{
        \cellcolor[rgb]{0.9,0.9,0.9}{
            \makecell[{{p{\linewidth}}}]{
                \texttt{\tiny{[GM$|$GM]}}
                ['clitoris', 'pyramid']
            }
        }
    }
    & & \\ \\

    \theutterance \stepcounter{utterance}  
    & & & \multicolumn{2}{p{0.3\linewidth}}{
        \cellcolor[rgb]{0.9,0.9,0.9}{
            \makecell[{{p{\linewidth}}}]{
                \texttt{\tiny{[GM$|$GM]}}
                {'word': 'clitoris', 'assignment': 'team'}
            }
        }
    }
    & & \\ \\

    \theutterance \stepcounter{utterance}  
    & & & \multicolumn{2}{p{0.3\linewidth}}{
        \cellcolor[rgb]{0.9,0.9,0.9}{
            \makecell[{{p{\linewidth}}}]{
                \texttt{\tiny{[GM$|$GM]}}
                {'word': 'pyramid', 'assignment': 'team'}
            }
        }
    }
    & & \\ \\

    \theutterance \stepcounter{utterance}  
    & & & \multicolumn{2}{p{0.3\linewidth}}{
        \cellcolor[rgb]{0.9,0.9,0.9}{
            \makecell[{{p{\linewidth}}}]{
                \texttt{\tiny{[GM$|$GM]}}
                \texttt{opponent has won} \\
            }
        }
    }
    & & \\ \\

    \theutterance \stepcounter{utterance}  
    & \multicolumn{4}{p{0.6\linewidth}}{
        \cellcolor[rgb]{0.95,0.95,0.95}{
            \makecell[{{p{\linewidth}}}]{
                \texttt{\tiny{[P2$\rightarrow$P1]}}
                \texttt{Post{-}episode Reflexion (not sent to GM):} \\
\texttt{New plan: My previous attempt failed because I only targeted two of the three words ("clitoris" and "pyramid") with the clue "TIP", leaving "storybook" unguessed and allowing the opponent to win. The strategy of grouping words is correct, but the execution must ensure *all* remaining team words are covered in a single turn when possible. Next time, I will analyze the full list of remaining words to find a single semantic link that connects all of them. For the set \{clitoris, pyramid, storybook\}, I should have used a clue like "END" (referring to the tip of the clitoris, the apex of the pyramid, and the end of a story) or "POINT" (if the partner can stretch the association to the point of a story). If a single word cannot cover all three, I must prioritize the most difficult word ("clitoris") with a clue that also covers at least one other, ensuring the remaining word is easy to guess in the subsequent turn, rather than leaving a word behind. The critical fix is to verify that the TARGETS list includes *every* remaining team word before submitting the clue.} \\
            }
        }
    }
    & & \\ \\

\end{supertabular}
}

\end{document}
